% Applications exponentielle et logarithme riemanniennes.
% Utilise dans la \cref{sec:context-riemann-notion}.
% Styles et briques communs aux figures de geometrie riemannienne du
% \cref{ch:context-riemann}. Ce fichier peut etre \input plusieurs fois.
\definecolor{rmanifold}{RGB}{90,90,90}
\definecolor{rtangent}{RGB}{50,110,120}
\definecolor{rgeodesic}{RGB}{40,70,140}
\definecolor{rvector}{RGB}{175,95,45}
\tikzset{
  riemann/.style={
    >={Stealth[length=4.5pt,width=3.5pt]},
    line join=round,
    font=\footnotesize,
  },
  % La nappe qui represente la variete, centree sur l'origine.
  rsurface/.pic={
    \fill[rmanifold!7]
      (-3.2,0.15) .. controls (-2.95,1.75) and (-1.5,2.35) .. (0,2.35)
      .. controls (1.5,2.35) and (2.95,1.7) .. (3.2,0.15)
      .. controls (1.9,-0.7) and (-1.9,-0.7) .. (-3.2,0.15) -- cycle;
    \draw[rmanifold!70, thick]
      (-3.2,0.15) .. controls (-2.95,1.75) and (-1.5,2.35) .. (0,2.35)
      .. controls (1.5,2.35) and (2.95,1.7) .. (3.2,0.15);
    \draw[rmanifold!70, thick]
      (-3.2,0.15) .. controls (-1.9,-0.7) and (1.9,-0.7) .. (3.2,0.15);
  },
  % Le plan tangent, a placer en un point donne via un shift.
  rplane/.pic={
    \fill[rtangent!10, draw=rtangent!65]
      (-1.35,-0.42) -- (1.55,-0.88) -- (2.05,0.22) -- (-0.85,0.68) -- cycle;
  },
}

\begin{tikzpicture}[riemann]
  \coordinate (P) at (-1.55,1.20);
  \coordinate (Q) at ( 1.75,1.45);

  \path (0,0) pic {rsurface};
  \node[rmanifold!85] at (2.75,0.75) {$\Manifold$};

  \path (P) pic[scale=0.8] {rplane};
  \node[rtangent!90, anchor=east, inner sep=1pt] at ($(P)+(-1.15,0.40)$)
    {$\Tangent{\param_0}{\Manifold}$};

  % la geodesique
  \draw[rgeodesic, very thick]
    (P) .. controls (-0.55,2.10) and (0.75,2.20) .. (Q);

  % le vecteur tangent
  \draw[rvector, very thick, ->] (P) -- ($(P)+(1.60,-0.23)$);
  \node[rvector, anchor=north, inner sep=2pt] at ($(P)+(0.80,-0.34)$)
    {$\tvec=\log_{\param_0}(\param_1)$};

  % le retour sur la variete
  \draw[rvector!60, dotted, very thick, ->]
    ($(P)+(1.63,-0.21)$) .. controls (0.75,0.55) and (1.60,0.70) .. ($(Q)+(0.02,-0.12)$);
  \node[rvector!85, anchor=north, inner sep=2pt, font=\scriptsize] at (1.05,0.62)
    {$\exp_{\param_0}(\tvec)$};

  \fill (P) circle (1.7pt);
  \node[anchor=east, inner sep=3pt] at (P) {$\param_0$};
  \fill (Q) circle (1.7pt);
  \node[anchor=west, inner sep=3pt] at (Q) {$\param_1$};
\end{tikzpicture}
